\section{Discussion}
\label{sec:discussion}

\textbf{Report the composition, not the model.} Any claim of the form ``model \emph{M} does not
use tools'' is, as measured, a claim about
\texttt{M $\times$ protocol $\times$ serialization $\times$ parser $\times$ stack}. Agent
evaluations should report at minimum (i) server-parsed calls, (ii) emitted-but-unparsed calls
under a stated criterion, (iii) request errors, and (iv) the exact serving configuration.
Reporting only (i) is the practice this paper argues against, and (ii) is where the 21$\times$
scale trend lives.

\textbf{Pre-flight, don't post-hoc.} A missing \texttt{-\/-enable-auto-tool-choice} is the only
failure that announces itself, and template inspection alone is insufficient
(\S\ref{sec:families}). We therefore treat a \texttt{tools}-bearing request returning HTTP 200 ---
\emph{and} a server-parsed call under \texttt{tool\_choice: required} --- as a required pre-flight
for any FC experiment, shipped as a 98-line script (\texttt{preflight\_toolcall.py}) that issues
one canonical request, asserts \texttt{tool\_calls} is non-empty with the expected \texttt{name}
and parseable \texttt{arguments}, then repeats under \texttt{required} as a positive control. It
runs in seconds and would have caught every silent failure in this paper.

\textbf{The guidance for this flag lives in the wrong place.} We surveyed the HuggingFace model
cards of nine families for any mention of \texttt{-\/-tool-call-parser}. Of the eight reachable
cards, none mentions it, while six recommend serving the model with vLLM. The publisher tells you
which server to use, the server asks you which parser to use, and neither treats the pairing as
its responsibility. We do \emph{not} read this as evidence that mismatch is common in deployment
--- we surveyed documents, not deployments, and vLLM ships dedicated parsers for some forty
families, so following \emph{its} table generally works. The claim is narrower: \textbf{the one
piece of configuration that can silently zero a capability measurement is documented by neither
party at the point of use.}

\textbf{Interface repair precedes protocol comparison.} On Llama the ReAct-vs-FC gap was 31 points
before the \texttt{strict} fix and 19 after, both components separately significant. On Qwen,
conditioning on parsed items, we do not detect a residual (\emph{p}=0.118) --- but that interval is
wide enough to contain effects as large as the unconditioned gap, so it is evidence of
\emph{insufficient resolution}, not equivalence. \textbf{Comparisons of agent protocols that do not
first exhaust the serving configuration matrix are not measuring protocols.}

\textbf{A cold-start hazard for agentic RL.} In the formal control, the mismatched parser admits
only 10 executions from 8{,}689 tight-classified call-like emissions while a tool-free reward path remains open;
the repaired parser admits 16{,}844 executed tool interactions. This isolates access to sampled tool
experience from learning from that experience: the first changes by more than three orders of magnitude,
while held-out rescues do not change. Practitioners should therefore verify both a
\textbf{non-trivial tool-execution count in the first training step} and a held-out measure of the
behaviour the tool is meant to teach. A rising reward curve establishes neither by itself.
